\documentclass[12pt,oneside,a4paper,parskip=half]{scrbook}

% --- Language & Basic Config ---
\usepackage[utf8]{inputenc}
\usepackage[T1]{fontenc}
\usepackage[ngerman,english]{babel}
\usepackage{lmodern}
\usepackage{microtype}
\usepackage{csquotes}
\usepackage{hyphenat}
\usepackage{newunicodechar}

% --- Layout ---
\usepackage[a4paper,left=20mm,right=20mm,top=20mm,bottom=25mm]{geometry}
\usepackage{setspace}
\onehalfspacing
\sloppy

% --- Graphics & Plots ---
\usepackage{pgfplots}
\pgfplotsset{compat=1.17}
\usepackage{tikz}
\usetikzlibrary{arrows.meta,positioning,calc,fit,backgrounds,shapes.geometric,shadows,patterns}

% --- Math & Tables ---
\usepackage{amsmath,amsfonts,amssymb}
\usepackage{graphicx}
\usepackage{float}
\usepackage{booktabs}
\usepackage{longtable}
\usepackage{tabularx}
\usepackage{placeins}
\usepackage{colortbl} % For table highlighting
\usepackage{caption}
\captionsetup[table]{skip=10pt}

% --- Listings ---
\usepackage{xcolor}
\usepackage{listings}
\colorlet{punct}{red!60!black}
\definecolor{background}{HTML}{F8F8F8}
\definecolor{delim}{RGB}{20,105,176}
\colorlet{numb}{magenta!60!black}
\definecolor{pblue}{rgb}{0.13,0.13,1}
\definecolor{pgreen}{rgb}{0,0.5,0}
\definecolor{pred}{rgb}{0.9,0,0}
\definecolor{pgrey}{rgb}{0.46,0.45,0.48}

\lstdefinestyle{code}{
    basicstyle=\ttfamily\footnotesize,
    columns=fullflexible,
    showstringspaces=false,
    numbers=left,
    numberstyle=\tiny\color{pgrey},
    stepnumber=1,
    numbersep=8pt,
    backgroundcolor=\color{background},
    commentstyle=\color{pgreen}\itshape,
    keywordstyle=\color{pblue}\bfseries,
    stringstyle=\color{pred},
    breaklines=true,
    breakatwhitespace=true,
    tabsize=4,
    keepspaces=true,
    frame=l,
    framesep=4pt,
    framerule=0pt,
    linewidth=\textwidth
}

\lstset{style=code}

% --- Metadata Definitions ---
\def\BaAuthor{Noah Raupold, Mohammad Ayoubi,\\ Jonathan Berwind}
\def\BaAuthorPDF{Noah Raupold, Mohammad Ayoubi, Jonathan Berwind}
\def\BaStudyProgram{Computer Science}
\def\BaType{AI \& Security Project III}
\def\BaTitle{Next-Gen Network Intrusion Detection}
\def\BaSubtitle{Self-Supervised Vector Space Engineering with DINOv3 \& SwiGLU}
\def\BaDeadline{\today}

\newcommand{\TitleGraphic}{%
    \IfFileExists{Logo.png}{%
        \includegraphics[width=0.5\textwidth]{Logo.png}%
    }{%
        \fbox{\parbox[c][3cm][c]{0.5\textwidth}{Logo.png fehlt}}%
    }%
}

\newcommand*{\forcetwosidetitle}{
    \begingroup
    \cleardoubleoddpage
    \KOMAoptions{titlepage=true}
    \csname @twosidetrue\endcsname
    \maketitle
    \endgroup
}

% --- Hyperlinks ---
\usepackage{hyperref}
\hypersetup{
    colorlinks=true,
    linkcolor=black,
    filecolor=magenta,
    urlcolor=cyan,
    pdfauthor={\BaAuthorPDF},
    pdftitle={\BaTitle},
    pdfsubject={Network Security}
}

\begin{document}

% --- Front Matter ---
    \selectlanguage{english}
    \frontmatter
    \titlehead{Technische Hochschule Würzburg-Schweinfurt (THWS)\\Faculty of Computer Science and Business Information Systems}
    \subject{\BaType}
    \title{\texorpdfstring{\BaTitle\\[5mm]\normalsize{\BaSubtitle}\\[15mm]\TitleGraphic}{\BaTitle}}
    \author{\BaAuthor}
    \date{Submitted: \BaDeadline}
    \forcetwosidetitle

    \tableofcontents

% --- Main Matter ---
    \mainmatter

    \chapter{Introduction}

    \section{The Cyber Security Challenge}
    Modern network infrastructures face an evolving landscape of cyber threats. Traditional Intrusion Detection Systems (IDS) often rely on signature-based methods or shallow machine learning models (like Logistic Regression). Our analysis of the \textbf{NSL-KDD} dataset identified critical failures in these baselines:
    \begin{enumerate}
        \item \textbf{Data Problem:} Extreme Class Imbalance. Normal and DoS traffic comprise $>98\%$ of the data, while critical attacks like \textit{User-to-Root (U2R)} and \textit{Remote-to-Local (R2L)} represent $<0.5\%$.
        \item \textbf{Statistical Bias:} Standard models minimize global loss by favoring majority classes, leading to high accuracy on "Normal" traffic but failing to detect critical intrusions (False Negatives).
    \end{enumerate}

    \section{Project Objective}
    This project introduces a **"Supervised Vector Space Engineering"** pipeline. Instead of simple classification, we warp the feature space using **Self-Supervised Learning (DINOv3)** and **Gram Matrix Analysis** to maximize the separation between normal traffic and rare anomalies. The final classification is performed by a Deep Neural Network enhanced with **SwiGLU activations**.

    \chapter{Architecture: The Pipeline}

    Our solution transitions from linear models to Deep Learning with Attention Mechanisms. The architecture consists of a robust preprocessing pipeline followed by a Teacher-Student training loop.

    \section{Data \& Preprocessing Pipeline}
    Before training, raw network data undergoes state-of-the-art (SOTA) preprocessing to handle skewness and feature engineering.

% --- TIKZ: PREPROCESSING PIPELINE ---
    \begin{figure}[H]
        \centering
        \begin{tikzpicture}[
            node distance=1.0cm,
            block/.style={rectangle, draw, fill=blue!10, text width=12em, text centered, rounded corners, minimum height=3em, drop shadow},
            arrow/.style={-Latex, thick}
        ]
            \node (raw) [block, fill=gray!20] {\textbf{Raw Data}\\(network\_connections.csv)};
            \node (split) [block, below=of raw] {Train/Test Split};

            \node (log) [block, below=of split] {\textbf{Log Transform}\\(Skewed Features)};
            \node (eng) [block, below=of log] {\textbf{Feature Engineering}\\(Byte Ratio, Srv Diff)};
            \node (scale) [block, below=of eng] {\textbf{StandardScaler}};
            \node (enc) [block, below=of scale] {\textbf{Label/Cat Encoding}};

            \node (gram) [block, below=of enc, fill=red!10] {\textbf{Gram Matrix Analysis}\\(Compute Feature Clusters)};

            \draw [arrow] (raw) -- (split);
            \draw [arrow] (split) -- (log);
            \draw [arrow] (log) -- (eng);
            \draw [arrow] (eng) -- (scale);
            \draw [arrow] (scale) -- (enc);
            \draw [arrow] (enc) -- (gram);

            % Grouping box
            \node[draw, dashed, inner sep=0.5cm, fit=(log) (gram), label=left:\rotatebox{90}{\textbf{SOTA Preprocessing}}] {};

        \end{tikzpicture}
        \caption{The Data Preprocessing Pipeline utilizing Gram Matrix Analysis for feature clustering.}
    \end{figure}

    \section{DINOv3 & Gram Anchoring Strategy}
    To learn robust features from imbalanced data, we implemented a Teacher-Student architecture (DINO) adapted for tabular data.

    \textbf{Core Components:}
    \begin{itemize}
        \item \textbf{Student Encoder:} Tabular ViT (Vision Transformer) receiving masked views of data.
        \item \textbf{Teacher Encoder:} Updated via Exponential Moving Average (EMA) of the student.
        \item \textbf{Loss Landscape:}
        \begin{enumerate}
            \item \textbf{DINO Loss:} Global alignment (Sinkhorn Knopp centering).
            \item \textbf{MIM Loss:} Masked Image Modeling (Reconstruction of masked tokens).
            \item \textbf{KoLeo Loss:} Enforces feature diversity in the batch.
            \item \textbf{Gram Matrix Loss:} Preserves structural correlations between features.
        \end{enumerate}
    \end{itemize}

% --- TIKZ: DINO ARCHITECTURE ---
    \begin{figure}[H]
        \centering
        \begin{adjustbox}{max width=\textwidth}
            \begin{tikzpicture}[
                node distance=1.5cm,
                auto,
                block/.style={rectangle, draw, fill=white, text width=7em, text centered, rounded corners, minimum height=3em, drop shadow},
                line/.style={draw, -Latex, thick},
                teacher/.style={draw, dashed, inner sep=0.3cm, rounded corners, fill=yellow!10},
                student/.style={draw, dashed, inner sep=0.3cm, rounded corners, fill=blue!10}
            ]

                % Input
                \node [block, fill=gray!20] (input) {Preprocessing Output};
                \node [block, below left=1.5cm and 1cm of input, fill=red!10] (mask) {Gram Masking\\(Noise Injection)};

                % Student Branch
                \node [block, below=1cm of mask, fill=blue!20] (s_enc) {Student Encoder\\(Tabular ViT)};
                \node [block, below=1cm of s_enc] (s_head) {Heads:\\DINO + MIM};

                % Teacher Branch
                \node [block, right=4cm of s_enc, fill=orange!20] (t_enc) {Teacher Encoder\\(EMA Update)};
                \node [block, below=1cm of t_enc] (t_head) {Head:\\DINO};

                % Losses
                \node [circle, draw, fill=white, below=2cm of s_head] (loss) {$\mathcal{L}$};

                % Connections
                \path [line] (input) -- (mask);
                \path [line] (mask) -- (s_enc);
                \path [line] (input) -| node[near start] {Unmasked View} (t_enc);

                \path [line] (s_enc) -- (s_head);
                \path [line] (t_enc) -- (t_head);

                \path [line] (s_head) -- node[left] {Prediction} (loss);
                \path [line] (t_head) |- node[near start] {Target} (loss);

                % EMA Arrow
                \draw [->, thick, dashed] (s_enc) -- node[above] {EMA Weights} (t_enc);

                % Loss Annotations
                \node [right=0.5cm of loss, text width=6cm, font=\small, align=left] (losstext) {
                    \textbf{Optimization Objectives:}\\
                    1. DINO Loss (Cross-Entropy)\\
                    2. MIM Loss (Reconstruction)\\
                    3. \textbf{KoLeo Loss} (Feature Diversity)\\
                    4. \textbf{Gram Matrix Loss} (Correlation)
                };

                % Backgrounds
                \begin{scope}[on background layer]
                    \node[student, fit=(mask) (s_enc) (s_head), label=above:\textbf{Student Branch}] {};
                    \node[teacher, fit=(t_enc) (t_head), label=above:\textbf{Teacher Branch}] {};
                \end{scope}

            \end{tikzpicture}
        \end{adjustbox}
        \caption{The Self-Supervised Training Architecture including KoLeo and Gram losses.}
        \label{fig:dino_archi}
    \end{figure}

    \section{Fine-Tuning: SwiGLU Classification Head}
    For the final classification, we replace standard ReLU layers with **SwiGLU** (Swish-Gated Linear Unit). This gating mechanism filters noise in the latent representation before the final Softmax layer, improving gradient flow.

    \chapter{Implementation Details}

    The implementation utilizes PyTorch. The core model relies on a custom `DINOLoss` for self-distillation and a `SwiGLU` activation for the classifier.

    \begin{lstlisting}[language=Python, caption=Dinov3 Loss & Architecture Implementation]
import torch
import torch.nn as nn
import torch.nn.functional as F

class DINOLoss(nn.Module):
    def __init__(self, out_dim, teacher_temp=0.04, student_temp=0.1, center_momentum=0.9):
        super().__init__()
        self.student_temp = student_temp
        self.teacher_temp = teacher_temp
        self.center_momentum = center_momentum
        self.register_buffer("center", torch.zeros(1, out_dim))

    def forward(self, student_output, teacher_output):
        # Teacher sharpening (Sinkhorn-Knopp style)
        t_out = F.softmax((teacher_output - self.center) / self.teacher_temp, dim=-1)
        t_out = t_out.detach() # Stop gradient to teacher

        # Student log-softmax
        s_out = F.log_softmax(student_output / self.student_temp, dim=-1)

        # Update center
        self.update_center(teacher_output)

        return -torch.sum(t_out * s_out, dim=-1).mean()

    @torch.no_grad()
    def update_center(self, teacher_output):
        batch_center = torch.sum(teacher_output, dim=0, keepdim=True)
        batch_center = batch_center / len(teacher_output)
        self.center = self.center * self.center_momentum + batch_center * (1 - self.center_momentum)

class SwiGLU(nn.Module):
    """
    Swish-Gated Linear Unit
    Formula: SwiGLU(x) = (xW + b) * Swish(xV + c)
    """
    def __init__(self, dim_in, dim_out):
        super().__init__()
        self.w = nn.Linear(dim_in, dim_out)
        self.v = nn.Linear(dim_in, dim_out)

    def forward(self, x):
        return self.w(x) * F.silu(self.v(x))
    \end{lstlisting}

    \chapter{Evaluation & Results}

    The model was evaluated on the **NSL-KDD Test+** dataset. We specifically focused on the F1-Score for the minority classes (U2R, R2L).

    \section{Binary Performance (Attack vs Normal)}
    The model demonstrates high capability in distinguishing general attacks from normal traffic.

    \begin{table}[H]
        \centering
        \renewcommand{\arraystretch}{1.2}
        \begin{tabular}{lcccc}
            \toprule
            \textbf{Class} & \textbf{Precision} & \textbf{Recall} & \textbf{F1-Score} & \textbf{Support} \\
            \midrule
            Normal & 0.8048 & \textbf{0.9382} & 0.8664 & 9711 \\
            Attack & \textbf{0.9465} & 0.8278 & \textbf{0.8832} & 12833 \\
            \midrule
            \textbf{Accuracy} & & & \textbf{87.54\%} & 22544 \\
            \textbf{Weighted Avg} & 0.8855 & 0.8754 & 0.8760 & 22544 \\
            \bottomrule
        \end{tabular}
        \caption{Binary Evaluation Metrics. The model achieves 87.54\% global accuracy on the binary task.}
    \end{table}

    \section{Multiclass Performance}
    While the baseline Logistic Regression model failed to generalize on sparse data (Recall < 50\% on U2R), our Tabular ViT architecture maintains detection capabilities across all classes.

    \begin{itemize}
        \item \textbf{Matthews Correlation Coefficient (MCC):} 0.7578
        \item \textbf{Global Accuracy:} 83.53\%
    \end{itemize}

    \begin{table}[H]
        \centering
        \renewcommand{\arraystretch}{1.2}
        \begin{tabular}{lcccc}
            \toprule
            \textbf{Class} & \textbf{Precision} & \textbf{Recall} & \textbf{F1-Score} & \textbf{Support} \\
            \midrule
            DoS & \textbf{0.9616} & 0.8433 & \textbf{0.8986} & 7458 \\
            Normal & 0.8048 & \textbf{0.9382} & 0.8664 & 9711 \\
            Probe & 0.8080 & 0.9124 & 0.8570 & 2421 \\
            \rowcolor{yellow!10} R2L & 0.6588 & 0.4080 & \textbf{0.5040} & 2887 \\
            \rowcolor{red!10} U2R & 0.2795 & 0.6716 & \textbf{0.3947} & 67 \\
            \midrule
            \textbf{Weighted Avg} & 0.8368 & 0.8353 & 0.8282 & 22544 \\
            \bottomrule
        \end{tabular}
        \caption{Final Multiclass Evaluation. Highlights include the non-zero F1-score for U2R (0.39) and R2L (0.50), significantly outperforming baselines which often yield 0.0.}
        \label{tab:results}
    \end{table}

    \section{Visualizing Performance}
    The confusion matrix highlights the model's ability to distinguish between normal traffic and attacks.

% --- TIKZ: METRICS VISUALIZATION ---
    \begin{figure}[H]
        \centering
        \begin{tikzpicture}
            \begin{axis}[
                ybar,
                bar width=0.8cm,
                symbolic x coords={DoS, Normal, Probe, R2L, U2R},
                xtick=data,
                ylabel={F1-Score},
                ymin=0, ymax=1,
                nodes near coords,
                nodes near coords style={font=\footnotesize},
                title={Per-Class Performance (F1-Score)},
                width=\textwidth,
                height=7cm,
                grid=major,
                ymajorgrids=true,
                axis on top
            ]
                \addplot[fill=blue!60, draw=black] coordinates {
                    (DoS,0.8986)
                    (Normal,0.8664)
                    (Probe,0.8570)
                    (R2L,0.5040)
                    (U2R,0.3947)
                };
            \end{axis}
        \end{tikzpicture}
        \caption{F1-Score distribution across classes. The model successfully detects rare attack vectors (R2L, U2R) that are usually missed.}
    \end{figure}

    \chapter{Conclusion}
    This portfolio demonstrated that **Supervised Vector Space Engineering**—combining DINO-based self-supervision with Gram Anchoring—effectively mitigates the class imbalance problem.

    \textbf{Key Findings:}
    \begin{enumerate}
        \item \textbf{Robustness:} Successfully implemented a State-of-the-Art Transformer architecture for NIDS.
        \item \textbf{SSL Superiority:} Self-Supervised Learning (SSL) significantly outperforms baselines in high-imbalance scenarios, achieving a multiclass accuracy of \textbf{83.53\%} and detecting 67\% of U2R attacks (Recall), whereas standard models often achieve 0\%.
    \end{enumerate}

\end{document}